\section{Electric Charge as Phase Asymmetry in Closure}
\label{sec:charge}

\subsection{Reformulating the Origin of Charge}

In the Standard Model, electric charge is assigned as a fundamental property
of particles.
Within the present framework, charge is examined as a structural feature
of triangular closure:
\begin{equation}
  \boxed{
    \text{Can electric charge be understood as an internal asymmetry
    of closure structure?}
  }
\end{equation}

\subsection{Symmetric Phase Configuration}

Consider a non-degenerate triangular closure.
Let each node carry a phase-like quantity $\phi_i$, normalized such that:
\begin{equation}
  \sum_{i=1}^{3} \phi_i = 1.
\end{equation}
The symmetric configuration is:
\begin{equation}
  \phi_1 = \phi_2 = \phi_3 = \frac{1}{3}.
\end{equation}
This configuration is structurally balanced --- no node is preferred,
no directional bias exists.
\begin{equation}
  \boxed{\text{The symmetric configuration corresponds to zero net charge.}}
\end{equation}

\subsection{Introducing Asymmetry}

When the phase distribution becomes asymmetric,
deviations from the mean arise.
Define charge-like quantities as:
\begin{equation}
  q_i = \phi_i - \frac{1}{3}.
\end{equation}
For example, a maximally asymmetric configuration
$(\phi_1, \phi_2, \phi_3) = (1, 0, 0)$ yields:
\begin{equation}
  q = \left(+\frac{2}{3},\; -\frac{1}{3},\; -\frac{1}{3}\right).
\end{equation}
\begin{equation}
  \boxed{
    \text{Fractional values such as } +\tfrac{2}{3} \text{ and } -\tfrac{1}{3}
    \text{ emerge naturally from three-node asymmetry.}
  }
\end{equation}

\subsection{Charge as Deviation from Symmetry}

\begin{equation}
  \boxed{
    \text{Charge can be interpreted as deviation from the symmetric
    phase configuration.}
  }
\end{equation}
This implies that charge is relational rather than intrinsic ---
it reflects internal imbalance within closure.

\subsection{Charge Conservation}

By construction:
\begin{equation}
  \sum_i q_i = 0.
\end{equation}
\begin{equation}
  \boxed{
    \text{Charge conservation follows from normalization of the
    internal structure.}
  }
\end{equation}
No additional conservation law needs to be imposed.

\subsection{Fractional Structure}

The appearance of fractional values is a direct consequence
of the three-node structure.
The natural unit is $1/3$.
\begin{equation}
  \boxed{
    \text{Fractional charges can be understood as a geometric consequence
    of triangular closure.}
  }
\end{equation}

\subsection{Degenerate Closure and Integer Charges}

In the degenerate regime, effective dimensionality reduces and
the three-node structure collapses toward a two-state system.
The phase structure simplifies, and integer-like charge values emerge.
\begin{equation}
  \boxed{
    \text{Integer charges can be associated with reduced (degenerate)
    closure structures.}
  }
\end{equation}
In the extreme degenerate limit, $q \approx 0$,
corresponding to minimal interaction.

\subsection{Structural Interpretation}

\begin{equation}
  \boxed{
    \text{Charge is not an assigned quantity, but a manifestation of
    internal asymmetry in closure geometry.}
  }
\end{equation}
This provides a structural basis for both fractional charge values
and charge conservation.
